% ----------------------
\subsection{Section 67 (Circa 2 November 1831)}
% ----------------------
	\label{DC67}
	
	% -----
	\subsubsection{Relation to section 1}
	% -----
	
		Section 67 was given at the same conference during which JS received \hyperref[DC1]{section 1}.
	
	% -----
	\subsubsection{Historical Background}
	% -----
		\label{DC67-HB}
		
		% --
		\paragraph{JSP}
		% --
		
			``At a 1--2 November 1831 conference that discussed the publication of the Book of Commandments, JS `received' a `testimony of the witnesses to the book of the Lord's commandments.' Several elders signed the statement and, by so doing, testified that the Holy Ghost had `born record' that the revelations `are given by inspiration of God \& are profitable for all men \& are verily true.'%
				\footnote{%
					% \label{}%
					See their ``\hyperref[EMS-1831-JS-CIR-2-NOV-1831-TEST]{testimony}.''
					}
			At this same conference, JS dictated this revelation, which addressed the fact that some of the elders had not received such a spiritual conviction. That manifestation had not come, the revelation explained, because of `fears' in the hearts of the elders. The revelation also disclosed that some of the elders questioned the `imperfections' in the language JS used in the revelations, even though another revelation that emerged from the conference declared that the commandments were given to God's `Servents in their weakness after the manner of their Language.'%
				\footnote{%
					% \label{}%
					This is a reference to \hyperref[DC1-24]{DC 1:24}, which was also given at this conference.
					}
			To dispel the hesitation to testify about the source of the revelations, this November revelation challenged the elders to write a revelation `like unto' those dictated by JS, emphasizing that if the elders failed, they would be justified in testifying to the world of the divine origin of JS's revelations and condemned if they did not so testify.
			
			``The original manuscript of this revelation is not extant, but John Whitmer dated it 2 November when he copied it into Revelation Book 1, likely sometime before his departure for Missouri on 20 November 1831. Unlike the minutes for proceedings of the conference on 1 November, the minutes for 2 November do not specifically mention dictation of a revelation, leaving open the possibility that this revelation could be the one referred to in the 1 November minutes as coming after JS `asked the conference what testimony they were willing to attach to these commandments.' According to the minutes, several elders `arose and said that they were willing to testify to the world that they knew that they were of the Lord.' Thereafter, a revelation was `received relative to the same'---perhaps this revelation, despite its being dated 2 November. A later JS history supports this possibility, stating that after the dictation of the preface to the Book of Commandments on 1 November, `some conversation was had concerning Revelations and language.' The history then presents this revelation as following that conversation.6 Although this is a plausible scenario for the production of the revelation, it does not explain why Whitmer then dated the document 2 November. Whatever the case, it is clear that the revelation---addressed to `ye Elders of my Church who have assembelled yourselves together'---was dictated sometime during the 1--2 November conference.''
			
	% -----
	\subsubsection{Versions}
	% -----
		\label{DC67-V}
		
		See the \href{https://drive.google.com/file/d/1goAvNz8jS4R8slYfMG_db55Ty2z_vmgK/view?usp=sharing}{sections comparison} document.
	
		\begin{compactdesc}
			\item[JSP] \hyperref[EMS-1831-JS-CIR-2-NOV-1831]{Circa 2 November 1831}
			\item[BC] ---
			\item[DC 1835] Section 25, 151--152
			\item[TS] 5:8 (15 April 1844), 496
			\item[DC 1844] Section 25, 221--222
			\item[MS] 5:12 (May 1845), 185
		\end{compactdesc}

	% -----
	\subsubsection{Section 67:1} % *DC67-1 
	% -----
		\label{DC67-1}

		BEHOLD and hearken, O ye elders of my church, who have assembled yourselves together, whose prayers I have heard, and whose hearts I know, and whose desires have come up before me.

		\begin{framed}
		\scriptsize
			\begin{compactdesc}
				\item[JSP] Behold \& hearken oh ye Elders of my Church who have assembelled yourselves together whose prayers I have heard \& whose \sout{hearts} desires have come up before me
				% \item[BC] 
				% \item[]
			\end{compactdesc}
		\end{framed}

	% -----
	\subsubsection{Section 67:2} % *DC67-2 
	% -----
		\label{DC67-2}

		Behold and lo, mine eyes are upon you, and the heavens and the earth are in mine hands, and the riches of eternity%
			\footnote{%
				% \label{}%
				For ``riches of eternity,'' see \hyperref[ETERNITY-PHRASES]{Eternity}. There are no biblical uses of the phrase; it's a purely restoration thing.
				}
		are mine to give.

		\begin{framed}
		\scriptsize
			\begin{compactdesc}
				\item[JSP] behold \& Lo mine eyes are upon you \& the heavens \& the earth are in mine hands \& the riches of eternity are mine to give[.]
				% \item[BC] 
				% \item[]
			\end{compactdesc}
		\end{framed}

	% -----
	\subsubsection{Section 67:3} % *DC67-3 
	% -----
		\label{DC67-3}

		Ye endeavored to believe that ye should receive the blessing which was offered unto you; but behold, verily I say unto you there were fears in your hearts, and verily this is the reason that ye did not receive.%
			\footnote{%
				% \label{}%
				I guess the question I want to ask is: what happened as a result of the fears? The fears changed the way these people thought and acted, apparently.
				}

		\begin{framed}
		\scriptsize
			\begin{compactdesc}
				\item[JSP] ye endeavour to believe that \sout{you} ye should receive the blessing which was offered unto you but behold verily I say unto you there were fears in your hearts \& verily this is the reason that ye did not receive
				% \item[BC] 
				% \item[]
			\end{compactdesc}
		\end{framed}

	% -----
	\subsubsection{Section 67:4} % *DC67-4 
	% -----
		\label{DC67-4}

		And now I, the Lord, give unto you a testimony of the truth of these commandments which are lying before you.

		\begin{framed}
		\scriptsize
			\begin{compactdesc}
				\item[JSP] \& now I the Lord give unto you a testimony of the truth of those commandments which \sout{I} are lying before you
				% \item[BC] 
				% \item[]
			\end{compactdesc}
		\end{framed}

	% -----
	\subsubsection{Section 67:5} % *DC67-5 
	% -----
		\label{DC67-5}

		Your eyes have been upon my servant Joseph Smith, Jun., and his language you have known, and his imperfections you have known; and you have sought in your hearts knowledge that you might express beyond his language; this you also know.

		\begin{framed}
		\scriptsize
			\begin{compactdesc}
				\item[JSP] your eyes have been upon my Servent Joseph \& his language you have known \& his imperfections you have known \& you have sought in your hearts knowlege that you might express beyond his language this you also know
				% \item[BC] 
				% \item[]
			\end{compactdesc}
		\end{framed}

	% -----
	\subsubsection{Section 67:6} % *DC67-6 
	% -----
		\label{DC67-6}

		Now, seek ye out of the Book of Commandments, even the least that is among them, and appoint him that is the most wise among you;

		\begin{framed}
		\scriptsize
			\begin{compactdesc}
				\item[JSP] now seek ye out of the Book of commandments even the least that is among them \& appoint him that is the most wise among you
				% \item[BC] 
				% \item[]
			\end{compactdesc}
		\end{framed}

	% -----
	\subsubsection{Section 67:7} % *DC67-7 
	% -----
		\label{DC67-7}

		Or, if there be any among you that shall make one like unto it, then ye are justified in saying that ye do not know that they are true;

		\begin{framed}
		\scriptsize
			\begin{compactdesc}
				\item[JSP] or if there be any among you that shall make one like unto it then ye are Justified in saying that ye do not know that is true
				% \item[BC] 
				% \item[]
			\end{compactdesc}
		\end{framed}

	% -----
	\subsubsection{Section 67:8} % *DC67-8 
	% -----
		\label{DC67-8}

		But if ye cannot make one like unto it, ye are under condemnation if ye do not bear record that they are true.

		\begin{framed}
		\scriptsize
			\begin{compactdesc}
				\item[JSP] but if you cannot make one like unto it ye are under condemnation if \sout{you} ye do not bear that it is true
				% \item[BC] 
				% \item[]
			\end{compactdesc}
		\end{framed}

	% -----
	\subsubsection{Section 67:9} % *DC67-9 
	% -----
		\label{DC67-9}

		For ye know that there is no unrighteousness in them, and that which is righteous cometh down from above, from the Father of lights.%
			\footnote{%
				% \label{}%
				This whole verse is clearly a reference to \hyperref[JAMES-1-17]{James 1:17}; for ``Father of lights,'' see \hyperref[LIGHT-PHRASES]{Light}. There are only two occurrences of that phrase. According to the \hyperref[JST-TIMELINE]{JST timeline}, JS wouldn't have been working on James in the NT translation until around the end of February 1832 or so.
				}

		\begin{framed}
		\scriptsize
			\begin{compactdesc}
				\item[JSP] for ye know that there is no unrighteousness in it \& that which is righteous cometh down from above from the father of lights
				% \item[BC] 
				% \item[]
			\end{compactdesc}
		\end{framed}

	% -----
	\subsubsection{Section 67:10} % *DC67-10 
	% -----
		\label{DC67-10}

		And again, verily I say unto you that it is your privilege, and a promise I give unto you that have been ordained unto this ministry, that inasmuch as you strip yourselves from jealousies and fears, and humble yourselves before me, for ye are not sufficiently humble, the veil shall be rent and you shall see me and know that I am---not with the carnal neither natural mind,%
			\footnote{%
				% \label{}%
				Note that this ``mind'' isn't in the JSP version. For these phrases involving ``mind,'' see \hyperref[MIND-PHRASES]{Mind}. Probably the most important connection mentioned there is to \hyperref[ROMANS-8]{Romans 8}.
				}
		but with the spiritual.

		\begin{framed}
		\scriptsize
			\begin{compactdesc}
				\item[JSP] \& again verily I say unto you that it is your privilege \& a promise I give unto you that have been ordained unto the ministry that in as much as ye strip yourselves from Jealesies \& fears \& humble yourselves before me for ye are not sufficiently humble the veil shall \sout{not} be wrent \& you shall see me \& know that I am not with the carnal neither natural but with the spiritual
				\item[DC 1835] 3 And again, verily I say unto you, that it is your privilege, and a promise I give unto you that have been ordained unto this ministry, that inasmuch as you strip yourselves from jealousies and fears, and humble yourselves before me, for ye are not sufficiently humble, the vail shall be rent and you shall see me and know that I am; not with the carnal, neither natural mind, but with the spiritual; 
				% \item[]
			\end{compactdesc}
		\end{framed}

	% -----
	\subsubsection{Section 67:11} % *DC67-11 
	% -----
		\label{DC67-11}

		For no man has seen God at any time%
			\footnote{%
				% \label{}%
				Compare \hyperref[JOHN-1-18]{John 1:18} and \hyperref[1JOHN-4-12]{1 John 4:12}.
				}
		in the flesh, except quickened%
			\footnote{%
				% \label{}%
				Note that the ``quickened'' isn't in the JSP version. The place in the DC for ``quicken'' is \hyperref[DC88]{section 88}, so maybe they added this word in later, thinking of those kinds of things.
				} 
		by the Spirit of God.

		\begin{framed}
		\scriptsize
			\begin{compactdesc}
				\item[JSP] for no man hath seen God at any time in the flesh but by the Spirit of God
				\item[DC 1835] for no man has seen God at any time in the flesh, except quickened by the Spirit of God:
				% \item[]
			\end{compactdesc}
		\end{framed}

	% -----
	\subsubsection{Section 67:12} % *DC67-12 
	% -----
		\label{DC67-12}

		Neither can any natural man abide the presence of God, neither after%
			\footnote{%
				% \label{}%
				Here are the fifth and sixth defintions of ``after'' from the 1828:
					\begin{compactitem}
						\item[5] According to
						\item[6] According to the direction and influence of
					\end{compactitem}
				So that's what we're dealing with here, I think.
				}
		the carnal mind.

		\begin{framed}
		\scriptsize
			\begin{compactdesc}
				\item[JSP] neither can any natural man abide the presence of God neither after the carnal mind
				\item[DC 1835] neither can any natural man abide the presence of God; neither after the carnal mind;
				% \item[]
			\end{compactdesc}
		\end{framed}

	% -----
	\subsubsection{Section 67:13} % *DC67-13 
	% -----
		\label{DC67-13}

		Ye are not able to abide the presence of God now, neither the ministering of angels;%
			\footnote{%
				% \label{}%
				Could be a reference to the kingdoms after the resurrection; think of \hyperref[DC76-86]{DC 76:86--88}.
				}
		wherefore, continue in patience until ye are \hyperref[PERFECTION]{perfected}.

		\begin{framed}
		\scriptsize
			\begin{compactdesc}
				\item[JSP] ye are not \sout{to} able to abide the presence of God now neither the ministering of Angels wherefore continue in patience untill ye are perfected
				\item[DC 1835] ye are not able to abide the presence of God now, neither the ministering of angels: wherefore continue in patience until ye are perfected.
				% \item[]
			\end{compactdesc}
		\end{framed}

	% -----
	\subsubsection{Section 67:14} % *DC67-14 
	% -----
		\label{DC67-14}

		Let not your minds turn back; and when ye are worthy, in mine own due time, ye shall see and know that which was conferred upon you by the hands of my servant Joseph Smith, Jun.  Amen.

		\begin{framed}
		\scriptsize
			\begin{compactdesc}
				\item[JSP] let not your minds turn back \& when ye are worthy in mine own due time ye shall see \& know that which was confirmed <upon you> by the hands of my Ser[v]ant Joseph Amen
				\item[DC 1835] 4 Let not your minds turn back; and when ye are worthy, in mine own due time, ye shall see and know that which was conferred upon you by the hands of my servant Joseph Smith, jr. Amen.
				% \item[]
			\end{compactdesc}
		\end{framed}